\begin{abstract}
\thispagestyle{empty}

Выпускная квалификационная работа посвящена разработке веб-приложения на основе RAG для работы с документацией PostgreSQL с учётом версии. Объект исследования -- предоставление справочной и обучающей информации по документации PostgreSQL. Предмет исследования -- архитектурные, алгоритмические и интерфейсные решения веб-приложения. Цель работы состоит в разработке приложения, которое по вопросу пользователя, выбранной версии PostgreSQL и режиму ответа извлекает релевантные фрагменты корпуса, проверяет подтверждения, формирует ответ с источниками и сохраняет историю запросов.

Работа состоит из введения, четырёх глав, заключения, списка использованных источников и приложений. Общий объём пояснительной записки составляет 99 страниц. В работу включены 5 приложений на 20 листах. В работе приведены 4 рисунка, 19 таблиц и 16 листингов, из них в приложениях размещены 4 рисунка, 1 таблица и 9 листингов. Библиографический список содержит 58 источников: 35 печатных источников и 23 электронных ресурса. Количество иностранных источников составляет 51 наименование.

Во введении описаны актуальность, объект, предмет, цель, задачи, методы исследования и практическая значимость. Первая глава посвящена предметной области, структуре документации PostgreSQL, проблеме смешения версий, подходам к поиску и требованиям. Во второй главе описаны архитектура, сценарий обработки запроса, диаграммы, модель данных, подготовка корпуса, поиск с учётом версии, повторное ранжирование и проверка подтверждений. В третьей главе приведена реализация серверной части, базы данных, маршрутов API, индексации корпуса, обработки запроса и пользовательского интерфейса. В четвёртой главе описаны программа испытаний, автоматические тесты, сценарные проверки, результаты запуска тестов и сопоставление требований. В заключении представлены результаты работы.
\end{abstract}
